\subsection{QC-008: exact evaluation and scalar-rank certificates for three tail families}

\paragraph{Statement QC-008.}
The following matrices of explicitly constructed centered highest-weight
polynomials on genuine nine-by-nine determinants have full row rank.
All entries and all rank calculations use exact arithmetic modulo $101$.
\[
\begin{array}{c|r|r|r|r}
\text{tail weight}&\text{functions}&\text{sample points}&
 \text{rank}&\text{selected minor}\\\hline
(8,8,8)&804&830&804&70\\
(9,6,6)&1250&2060&1250&88\\
(9,7,4)&2030&2060&2030&9
\end{array}
\tag{QC-008.1}
\]
The first matrix was recomputed entry by entry independently of its
accelerated builder. The other two matrices were produced directly by
the standard-library evaluator. All three minors were then computed
with ordinary scalar Gaussian elimination, independently of the packed
arithmetic used for polynomial evaluation.

Combined with the complete highest-weight basis certificates and the
stabilization proofs S-005 and OBS-005, this proves
\[
 K_9(9d-24,8,8,8)=a_{(9d-24,8,8,8)},
\]
\[
 K_9(9d-21,9,6,6)=a_{(9d-21,9,6,6)},\qquad
 K_9(9d-20,9,7,4)=a_{(9d-20,9,7,4)}
\tag{QC-008.2}
\]
for every $d$ for which the displayed tuple is a partition.
In the first family the padded multiplicity also equals the ambient
multiplicity by S-005; in the other two, the universal inequality
$q\leq a=K_9$ suffices to exclude a multiplicity gap. No padded
multiplicity equality is inferred there.

\begin{proof}
The source samples are those of QC-007. The first matrix uses all
$830$ points with jets through order eight. The other matrices use
the extension to $2060$ points and jets through order nine, generated
by exactly the same coefficient interpolation and normalization.
All $6180$ additional direct determinant consistency evaluations passed.
The source forms have $F_0=1$ and $F_1=0$, so evaluating a centered
invariant or highest-weight polynomial is exactly the evaluation of
its denominator-cleared ambient representative at these points.

The integral coefficient arrays for the first basis use the color-orbit
encoding described in QC-007. For an orbit coefficient $q_O$, its
evaluation is $q_O$ times the sum of the six globally permuted
representative monomials, counted with repetition. This is equivalent
to giving every distinct orbit member coefficient $6q_O/|O|$.
The coefficient arrays for the other two bases already contain the
full monomials and need no orbit restoration. Their multidegrees,
highest-weight identities, and complete ranks as coefficient
polynomials are checked separately in the basis certificates.

Here is the arithmetic justification for evaluating every sample at
once in a standard-library implementation. Let $N$ be the number
of samples and encode a vector $a_i\in\{0,\ldots,100\}$ by
\[
 A=\sum_{i=0}^{N-1}a_i2^{32i}.
\]
For the values $b_i$ of a fixed tensor coordinate, prepare seven masks
\[
 M_k=\sum_{i:\,\text{bit}_k(b_i)=1}127\,2^{32i},
 \qquad k=0,\ldots,6.
\]
Then
\[
 Z=\sum_{k=0}^{6}(A\mathbin{\&}M_k)2^k
       =\sum_i a_ib_i2^{32i}.
 \tag{QC-008.3}
\]
Each lane is smaller than $101^2<2^{14}$, so there is no carry
between lanes. Products of monomials are obtained by repeating this
operation and reducing each lane after every multiplication. Caching
prefix products changes no arithmetic.

To verify the simultaneous reduction, put $r=\lfloor2^{14}/101\rfloor=162$.
For $0\leq z<101^2$, the number
$\lfloor zr/2^{14}\rfloor$ is either $\lfloor z/101\rfloor$ or
one less. The residual is nonnegative and below $202$, and subtracting
$101$ once if it is at least $101$ gives the correct residue. The
implementation performs this using a seven-bit mask for the quotient
and the bit of $z-101q+155$ at position eight for the final correction.
The intermediate product $zr$ occupies fewer than $21$ bits. Thus
$32$-bit lanes isolate the operations, and no subtraction borrows from
another lane. In addition to this proof, every one of the $10201$
possible inputs to the reduction and every one of the $10201$
pointwise products in $\mathbb F_{101}^2$ was checked exhaustively.

For a final linear combination, coefficients are reduced to
$0,\ldots,100$ first. An ordinary monomial evaluation is at most $100$
per lane, and a six-color orbit sum is at most $600$. The scripts
assert before each row that
\[
 100^2\,\#\text{monomials}<2^{32}
 \quad\text{or respectively}\quad
 6\cdot100^2\,\#\text{orbit monomials}<2^{32}.
 \tag{QC-008.4}
\]
Consequently the unreduced polynomial sums also have no inter-lane
carry. Unpacking and ordinary integer reduction modulo $101$ at the
end gives their exact values.

The program \texttt{verify\_centered\_evaluations.py} independently
reads the integral coefficient files, restores the color orbits, and
recomputed all $804\cdot830=667320$ entries of the first matrix.
Every entry agreed with the independently built matrix. The complete
check took $164.586$ seconds in this environment. The program
\texttt{evaluate\_full\_tail\_basis.py} uses the same proved arithmetic
for full monomials and writes all $1250\cdot2060$ and
$2030\cdot2060$ entries of the other two matrices directly. Profile
blocks are cached with both coefficient-file and sample-file hashes.
Neither program imports NumPy, SciPy, or a symbolic algebra system.

Finally \texttt{check\_centered\_rank.py} uses only ordinary scalar
Gaussian elimination on the stored bytes. For the first matrix it
returns rank $804$ and minor $70$ in $11.612$ seconds. For the second,
the first $1280$ columns already have rank $1250$ and minor $88$
($39.035$ seconds). For the third, all $2060$ columns give rank $2030$
and minor $9$ ($162.359$ seconds). The selected column indices are
recorded in the result files.

The matrix-byte SHA-256 digests, in function-major and then sample-major
order, are respectively
\begin{quote}\ttfamily\small
71b43abba722f968438ea9a4c77ad7f899f5a1dd6b8666962bdceabc8fd50d2d,\\
1af48de765acb9aac9b95a7e13d807589825028dbc32d563422a7714c295c4a2,\\
c7a31973e8a2c16b15323510cd8ebab81efe9682023a7f0bdfeaa071275955ce.
\end{quote}
These nonzero minors prove that every function in each complete
highest-weight basis survives restriction to the determinant variety
and that the restrictions are linearly independent. The samples have
integer source pencils; the tensor normalizations only have factorial
denominators prime to $101$. Thus a nonzero modular minor proves a
nonzero characteristic-zero minor. No modular semisimplicity assumption
is made. The all-degree consequences in (QC-008.2) now follow from the
actual-space stabilization and multiplication-by-$F_0$ argument of
S-005 and OBS-005.
\end{proof}

\paragraph{Scope.}
These are three families with fixed tails as $d$ varies; they are not
families obtained by multiplying every part of one partition by an
arbitrary stretch. For example, the calculation excludes
$(114,9,6,6)$ in degree $15$, but does not by itself decide a label
with a larger tail. In the octic-polar hypersurface terminology of
S-005 it also forces the generator's tail parameter $c_G$ to be at
least nine.

\paragraph{Signed nine-bracket extension of the evaluation checker.}
The program \texttt{verify\_centered\_evaluations\_c9.py} applies the same
exact arithmetic to the nine-bracket basis specialized at $C_2=I$.
It first verifies every quadratic tensor entry of every sample.
The quadratic vertices are then omitted from the monomial encoding;
the remaining tensor orders are three through nine.  Since nine is
odd, the restored global-color orbit is
\[
 q_O\sum_{\pi\in S_3}\operatorname{sgn}(\pi)\,
                    \pi\cdot\text{monomial}(O).
\]
Orbits with an odd stabilizer vanish and are absent from the archives.
Even stabilizers retain their required multiplicities.  The verifier
asserts this by collecting identical permuted monomials with signs
and requiring the sum of absolute signed multiplicities to be six.

For packed evaluation let $U_+$ and $U_-$ be the sums for the three
even and three odd permutations.  Every lane of either sum lies
between zero and $300$.  Hence
\[
               U_+ +303\sum_i2^{32i}-U_-
\]
has every lane between $3$ and $603$, has no inter-lane borrow,
and has the desired residue modulo $101$.  The already proved
reducer applies because $603<101^2$.  Reduce this signed orbit
before the final linear combination.  Consequently the ordinary
bound $100^2\#\text{orbit monomials}<2^{32}$ suffices for each row.
Each completed profile is checked against every saved sample entry
and cached with the SHA-256 digests of its coefficient archive,
sample file, and expected-value block.  Atomic certificate replacement
and hash matching allow an expanding basis to be verified without
treating a changed input as an already checked block.
This describes a checker and its correctness; a full nine-bracket
rank conclusion requires its separate completed result and scalar
rank certificate.

\paragraph{Completed nine-bracket arithmetic result.}
The final signed-orbit matrix has $2386$ rows and $2420$ columns.
The standard-library checker independently recomputed all
$5\,774\,120$ entries, with no discrepancies.  A separate scalar
Gaussian elimination, using no packed elimination arithmetic, found
rank $2386$ and determinant $14$ modulo $101$ on the first $2386$
columns.  Both computations bind the matrix to SHA-256
\begin{quote}\ttfamily\small
d781202425cd3e9bba02d996afa76de3b1bda036564bab797f3ac7135d439b50.
\end{quote}
The reports are \texttt{c9\_evaluation\_verification.json} and
\texttt{c9\_evaluation\_scalar\_rank.json} in this branch's result
directory.  The scalar computation took $383.717$ seconds.
The basis reconstruction and its completeness are separate checks
owned by the support branch; these arithmetic results provide the
evaluation and rank parts of that branch's fixed-tail theorem.
